\clearpage
\appendix

\section{Complete Experiment Log}
\label{app:experiments}

Table~\ref{tab:exp_log} lists all 26 experiments with full configurations and results. All experiments use 4-bit QLoRA via Unsloth on NVIDIA A100-SXM4-40GB GPUs. Evaluation uses greedy decoding (\texttt{do\_sample=False}) from Exp.~7 onwards; earlier experiments used stochastic sampling (marked with $\dagger$). Bold rows indicate best results per dataset size. ES = early stopping (patience=2); the notation ``8(ES$\to$4)'' means max 8 epochs with early stopping selecting epoch 4.

\begin{table*}[t]
\centering
\caption{Complete experiment log. All models are Qwen3-8B (4-bit) except Exp.~1 (Llama 3.1 8B). LoRA config is $r$/$\alpha$ with dropout 0.05 unless noted. $\dagger$ = stochastic eval (inflated). $\ddagger$ = thinking enabled (mismatch).}
\label{tab:exp_log}
\resizebox{\textwidth}{!}{%
\begin{tabular}{rlllrrccccccrr}
\toprule
Exp & Dataset & $N$ & Folds & Ep. & LR & NEF & Method & Baseline & FT & $\Delta$ & W$\to$R & R$\to$W & $p$ \\
\midrule
\multicolumn{14}{l}{\textit{Round 1: Model Selection}} \\
1   & cleaned (all)    & 256 & 5  & 3        & 2e-4   & --  & SFT$^\dagger$          & 37.5 & 39.8 & +2.4  & 35  & 29 & 0.25  \\
2   & cleaned (num)    & 215 & 5  & 4        & 5e-5   & --  & SFT$^{\dagger\ddagger}$ & 67.9 & 67.0 & $-$0.9 & 25  & 27 & 0.88  \\
\midrule
\multicolumn{14}{l}{\textit{Round 2: Hyperparameter Search (stochastic eval)}} \\
3   & cleaned (num)    & 215 & 5  & 3        & 2e-4   & 5   & SFT$^\dagger$  & 67.9 & 74.9 & +7.0$^\dagger$ & 31  & 16 & 0.13  \\
4   & cleaned (num)    & 215 & 5  & 3        & 2e-4   & 10  & SFT$^\dagger$  & 70.7 & 72.1 & +1.4  & 21  & 18 & 1.0   \\
5   & cleaned (num)    & 215 & 5  & 3        & 2e-4   & 5   & SFT$^\dagger$  & 67.9 & 69.3 & +1.4  & 24  & 21 & 0.88  \\
6   & cleaned (num)    & 215 & 5  & 3        & 1e-4   & 5   & SFT$^\dagger$  & 68.8 & 68.4 & $-$0.5 & 20  & 21 & 1.0   \\
\midrule
\multicolumn{14}{l}{\textit{Round 3: Deterministic Evaluation (greedy decoding)}} \\
7   & v2 (num)         & 280 & 5  & 3        & 2e-4   & 5   & SFT    & 64.0 & 65.2 & +1.3  & 21  & 21 & 0.88  \\
8   & v2 (num)         & 280 & 5  & 5        & 2e-4   & 5   & SFT    & 64.3 & 65.7 & +1.4  & 30  & 19 & 0.88  \\
9   & cleaned (num)    & 215 & 5  & 5        & 2e-4   & 5   & SFT$^\dagger$  & 70.7 & 72.6 & +1.9  & 21  & 17 & 0.50  \\
11  & cleaned (num)    & 215 & 5  & 3        & 2e-4   & 5   & SFT    & 70.7 & 71.6 & +0.9  & 22  & 20 & 0.88  \\
\textbf{12} & \textbf{cleaned (num)} & \textbf{215} & \textbf{5} & \textbf{4} & \textbf{2e-4} & \textbf{5} & \textbf{SFT} & \textbf{70.7} & \textbf{73.0} & \textbf{+2.3} & \textbf{23} & \textbf{18} & \textbf{0.50} \\
\textbf{16} & \textbf{v2 (num)} & \textbf{280} & \textbf{5} & \textbf{4} & \textbf{2e-4} & \textbf{5} & \textbf{SFT} & \textbf{71.0} & \textbf{73.0} & \textbf{+2.0} & \textbf{19} & \textbf{15} & \textbf{0.69} \\
17  & cleaned (num)    & 215 & 5  & 4        & 2e-4   & 7   & SFT    & 71.0 & 69.8 & $-$1.3 & 19  & 18 & 0.38  \\
\midrule
\multicolumn{14}{l}{\textit{Round 4: Alternative Methods}} \\
10  & v2 (num)         & 280 & 5  & DPO      & 5e-6   & --  & DPO    & 64.3 & --   & --    & --  & -- & --    \\
14  & cleaned (num)    & 215 & 5  & DPO      & 5e-6   & --  & DPO    & 69.8 & 68.0 & $-$1.7 & 4   & 7  & 0.63  \\
13  & cleaned (num)    & 215 & 5  & SFT+GRPO & 2e-4   & 5   & GRPO   & 70.7 & --   & --    & --  & -- & --    \\
15  & cleaned (num)    & 215 & 5  & SFT+GRPO & 2e-4   & 5   & GRPO   & 70.7 & 71.2 & +0.5  & 18  & 17 & 0.88  \\
\midrule
\multicolumn{14}{l}{\textit{Round 5: Paraphrase Augmentation (5-fold CV)}} \\
\textbf{18} & \textbf{v3 (num)} & \textbf{501} & \textbf{5} & \textbf{4} & \textbf{2e-4} & \textbf{5} & \textbf{SFT} & \textbf{59.1} & \textbf{65.3} & \textbf{+6.2} & \textbf{74} & \textbf{43} & \textbf{0.063} \\
19  & v3 (num)         & 501 & 5  & 8(ES$\to$4) & 1.5e-4 & 5 & SFT  & 59.1 & 64.5 & +5.4  & 83  & 56 & 0.063 \\
20  & v3 (num)         & 501 & 5  & 6        & 1e-4   & 5   & SFT    & 59.1 & 64.9 & +5.8  & 78  & 49 & 0.063 \\
\midrule
\multicolumn{14}{l}{\textit{Round 6: Scaled Augmentation (10-fold CV)}} \\
23  & v4 (num)         & 866 & 10 & 3        & 2e-4   & 5   & SFT    & 63.6 & 75.1 & +11.4 & 171 & 72 & \textbf{0.002} \\
21  & v4 (num)         & 866 & 10 & 4        & 2e-4   & 5   & SFT    & 63.6 & 79.3 & +15.7 & 194 & 58 & \textbf{0.002} \\
22  & v4 (num)         & 866 & 10 & 6(ES$\to$4) & 2e-4 & 5  & SFT    & 63.6 & 81.8 & +18.1 & 208 & 51 & \textbf{0.002} \\
\textbf{24} & \textbf{v4 (num)} & \textbf{866} & \textbf{10} & \textbf{8(ES$\to$4)} & \textbf{2e-4} & \textbf{5} & \textbf{SFT} & \textbf{63.6} & \textbf{82.2} & \textbf{+18.6} & \textbf{206} & \textbf{45} & \textbf{0.002} \\
\midrule
\multicolumn{14}{l}{\textit{Round 7: 600-Sample Subset Validation (10-fold CV)}} \\
\textbf{25} & \textbf{v4 sub (num)} & \textbf{600} & \textbf{10} & \textbf{5} & \textbf{2e-4} & \textbf{5} & \textbf{SFT} & \textbf{68.0} & \textbf{79.7} & \textbf{+11.7} & \textbf{108} & \textbf{38} & \textbf{0.002} \\
26  & v4 sub (num)     & 600 & 10 & 5(ES)    & 2e-4   & 5   & SFT    & 68.0 & 78.8 & +10.8 & 100 & 35 & \textbf{0.002} \\
\bottomrule
\end{tabular}%
}
\end{table*}

Table~\ref{tab:grpo_exp_log} lists all GRPO experiments. All use the custom DAPO-style training loop on 266 mixed-difficulty questions (unless noted), with Qwen3-8B 4-bit and LoRA $r=16$/$\alpha=32$. Evaluation uses deterministic seeding (seed $=3407+q_i$), temperature 1.0, top-$p$ 0.95.

\begin{table*}[t]
\centering
\caption{GRPO experiment log. All models are Qwen3-8B (4-bit) with 200 useful steps, $G=4$ generations, grad\_accum=4. Accuracy is on the 266 mixed questions unless noted.}
\label{tab:grpo_exp_log}
\resizebox{\textwidth}{!}{%
\begin{tabular}{lllccccccl}
\toprule
Name & Init & LR & Waste\% & Attempts & Step 80 & Step 100 & Step 200 & Holdout (54q) & Notes \\
\midrule
GRPO baseline        & SFT fold\_4 & 1e-5 & 23.1 & 260 & \textbf{57.9\%} & 54.1\% & 52.3\% & 53.7\% (step 80) & Best config \\
GRPO aggressive LR   & SFT fold\_4 & 2e-5 & 20.9 & 253 & \multicolumn{3}{c}{Evaluation pending} & --- & 2$\times$ LR \\
GRPO conservative LR & SFT fold\_4 & 5e-6 & 21.3 & 254 & \multicolumn{3}{c}{Evaluation pending} & --- & 0.5$\times$ LR \\
GRPO from scratch    & Fresh LoRA  & 1e-5 & 29.3 & 283 & 50.8\% & Pending & Pending & 55.6\% (step 200) & No SFT init \\
\bottomrule
\end{tabular}%
}
\end{table*}


Notes: Exp.~1 uses Llama 3.1 8B; all others use Qwen3-8B. Exp.~5 uses LoRA $r=8$/$\alpha=16$; all others use $r=16$/$\alpha=32$. Exp.~10 timed out after 14 hours. Exp.~13 crashed due to Triton kernel incompatibility. Exp.~14 completed only 4 of 5 folds (missing adapter for fold~4). Accuracy values are percentages; $\Delta$ is the absolute difference (FT $-$ Baseline). W$\to$R and R$\to$W are question flip counts aggregated across all folds.


\section{Key Prompts}
\label{app:prompts}

Below are selected prompts used in our data pipeline. Minor formatting adjustments have been made for readability.

\subsection{Synthetic Question Generation}
\label{app:prompt:gen}

Used to generate new reliability engineering problems from seed examples (Section~\ref{sec:data:selfinst}).

\begin{quote}
\small
\textit{You are an expert in reliability engineering, creating practice problems for students.}

\textit{I will show you \{num\_examples\} example problems. Your task is to:}
\begin{enumerate}\small
    \item \textit{Understand the style, difficulty, and topics of these examples}
    \item \textit{Create \{num\_to\_generate\} new, original problems that are similar in style and difficulty}
    \item \textit{Make questions that are educational and challenging}
    \item \textit{Provide complete solutions with step-by-step reasoning}
\end{enumerate}

\textit{Important: Create different questions (don't just change numbers in the examples). Use similar concepts but new scenarios. Include all necessary data in the question (self-contained). Show clear reasoning steps. Provide specific numerical answers when appropriate.}

\textit{[3 randomly sampled seed examples shown]}

\textit{For each problem, provide:}\\
\textit{QUESTION: [Complete problem statement]}\\
\textit{REASONING: [Step-by-step solution]}\\
\textit{ANSWER: [Final answer(s)]}
\end{quote}

\subsection{Paraphrase Augmentation}
\label{app:prompt:rephrase}

Used to rephrase existing questions while preserving mathematical content (Section~\ref{sec:data:augment}). The prompt includes two manually curated few-shot examples (omitted for brevity).

\begin{quote}
\small
\textit{You are a technical writer specializing in reliability engineering. Your task is to rephrase a question and its solution reasoning while preserving their exact mathematical meaning.}

\textit{Strict rules:}
\begin{enumerate}\small
    \item \textit{All numerical values, failure rates, probabilities, time periods, and parameters must remain exactly the same}
    \item \textit{All formulas, equations, and mathematical relationships must be preserved}
    \item \textit{The correct answer must be identical}
    \item \textit{Change the wording, sentence structure, and phrasing to create a distinct version}
    \item \textit{Do not add information, assumptions, or conditions not present in the original}
\end{enumerate}

\textit{What to change: Synonyms for non-technical words (e.g., ``compute'' $\to$ ``determine''), sentence structure, scenario framing (e.g., ``plant'' $\to$ ``factory''), transition words.}

\textit{What must not change: Any number, rate, probability, or parameter value. Variable names in formulas. The mathematical operations and formulas used. The final answer.}

\textit{[2 few-shot examples shown]}

\textit{Respond with a JSON object:}
\texttt{\{"question": "...", "reasoning": "..."\}}
\end{quote}

\subsection{Paraphrase Verification}
\label{app:prompt:verify}

Used by GPT-5.4 to verify that a rephrase preserves mathematical meaning (Section~\ref{sec:data:augment}).

\begin{quote}
\small
\textit{You are verifying whether a rephrased reliability engineering question preserves the exact mathematical meaning of the original. Compare the original and rephrased versions. Check all of the following:}
\begin{enumerate}\small
    \item \textit{Numerical values: Are all numbers, rates, probabilities exactly the same?}
    \item \textit{Mathematical meaning: Do both questions ask for the exact same quantity?}
    \item \textit{Constraints: Are all conditions and given information preserved?}
    \item \textit{Solution: Would the exact same calculation steps produce the correct answer?}
    \item \textit{Reasoning: Are all formulas and intermediate values identical?}
\end{enumerate}

\textit{Respond with a JSON object:}
\texttt{\{"values\_preserved": bool, "meaning\_preserved": bool, "constraints\_preserved": bool, "solution\_equivalent": bool, "reasoning\_equivalent": bool, "overall\_valid": bool, "issues": "..."\}}
\end{quote}

\subsection{GRPO System Prompt}
\label{app:prompt:grpo}

Used during GRPO training and evaluation (Section~\ref{sec:grpo}). The \texttt{/no\_think} prefix disables Qwen3's internal thinking mode for consistency with SFT format.

\begin{quote}
\small
\texttt{/no\_think}\\
\textit{You are a Reliability Engineering Expert. Solve the user's problem step-by-step with rigorous mathematical reasoning.}

\textit{Rules for your final answer:}
\begin{itemize}\small
    \item \textit{Write ONE single} \verb|\boxed{}| \textit{at the very end of your response --- your final answer only.}
    \item \textit{Do NOT use} \verb|\boxed{}| \textit{for intermediate steps or calculations.}
\end{itemize}
\textit{Always put your single final numerical answer inside} \verb|\boxed{}|\textit{.}
\end{quote}

\subsection{SFT System Prompt}
\label{app:prompt:system}

Used as the system message during both training and evaluation (Section~\ref{sec:sft:setup}).

\begin{quote}
\small
\textit{You are an expert in reliability engineering, probability theory, and system analysis. Provide clear, accurate answers with step-by-step reasoning. At the end, clearly state your final answer after `Final Answer:'.}
\end{quote}


\section{Project Organization}
\label{app:organization}

\subsection{Task Distribution}

This project was carried out by two students under the supervision of Prof.\ Zhiguo Zeng. Table~\ref{tab:tasks} summarizes the main tasks and their distribution.

\begin{table}[h]
\centering
\small
\begin{tabular}{@{}lcc@{}}
\toprule
\textbf{Task} & \textbf{Alexandre} & \textbf{Elora} \\
\midrule
Literature review          & \checkmark & \checkmark \\
Seed dataset curation      & \checkmark & \\
Synthetic data generation  & \checkmark & \\
Paraphrase augmentation    & \checkmark & \\
SFT training pipeline      & \checkmark & \\
GRPO implementation        & \checkmark & \checkmark \\
Evaluation \& baselines    & \checkmark & \checkmark \\
HPC cluster management     & \checkmark & \\
Paper writing (draft)      & \checkmark & \\
Paper writing (revision)   & \checkmark & \checkmark \\
Figures \& tables          &            & \checkmark \\
\bottomrule
\end{tabular}
\caption{Task distribution between team members.}
\label{tab:tasks}
\end{table}

\subsection{Project Timeline}

The project spanned approximately six months, from late October 2025 to April 2026. Figure~\ref{fig:gantt} presents the timeline as a Gantt chart.

\begin{figure*}[t]
\centering
\resizebox{\textwidth}{!}{%
\begin{ganttchart}[
    hgrid,
    vgrid={*{3}{draw=none}, dotted},
    x unit=0.45cm,
    y unit chart=0.6cm,
    title height=1,
    bar height=0.6,
    bar label font=\small,
    title label font=\small,
    milestone label font=\small,
    group label font=\small\bfseries,
    bar/.append style={fill=cvprblue!60},
    group/.append style={fill=cvprblue!30},
    milestone/.append style={fill=red!70},
]{1}{28}
    % Title: 7 months, 4 weeks each
    \gantttitle{Oct 25}{4}
    \gantttitle{Nov 25}{4}
    \gantttitle{Dec 25}{4}
    \gantttitle{Jan 26}{4}
    \gantttitle{Feb 26}{4}
    \gantttitle{Mar 26}{4}
    \gantttitle{Apr 26}{4} \\

    % Phase 1: Data
    \ganttgroup{Data Pipeline}{1}{20} \\
    \ganttbar{Literature review}{1}{8} \\
    \ganttbar{Seed dataset curation}{5}{14} \\
    \ganttbar{Synthetic data generation}{12}{18} \\
    \ganttbar{Paraphrase augmentation}{20}{22} \\

    % Phase 2: Training
    \ganttgroup{Model Training}{18}{26} \\
    \ganttbar{SFT pipeline \& experiments}{18}{23} \\
    \ganttbar{GRPO implementation}{20}{24} \\
    \ganttbar{Evaluation \& baselines}{22}{26} \\

    % Phase 3: Paper
    \ganttgroup{Paper}{24}{28} \\
    \ganttbar{Writing \& revision}{24}{28} \\

    % Milestones
    \ganttmilestone{Mid-project defence}{7} \\
    \ganttmilestone{Seed dataset finalized}{14} \\
    \ganttmilestone{Best SFT result (+18.6\%)}{23} \\
    \ganttmilestone{Submission}{28}

    % Links
    \ganttlink{elem2}{elem3}
    \ganttlink{elem3}{elem4}
    \ganttlink{elem6}{elem7}
\end{ganttchart}
}
\caption{Project timeline (October 2025 -- April 2026).}
\label{fig:gantt}
\end{figure*}

\subsection{Time Estimates}

Table~\ref{tab:time} provides approximate time spent on each major phase of the project.

\begin{table}[h]
\centering
\small
\begin{tabular}{@{}lr@{}}
\toprule
\textbf{Phase} & \textbf{Approx.\ hours} \\
\midrule
Literature review \& problem formulation   & 30 \\
Data pipeline (seed + synthetic + augment.) & 60 \\
SFT training \& hyperparameter search       & 50 \\
GRPO implementation \& experiments          & 40 \\
Evaluation \& baseline comparison           & 25 \\
Paper writing \& revision                   & 45 \\
\midrule
\textbf{Total (combined)}                   & \textbf{250} \\
\bottomrule
\end{tabular}
\caption{Approximate time spent on each project phase (combined for both team members).}
\label{tab:time}
\end{table}
